% =====================================================================
% Слайд 5 — Данные: исполнение в реальном времени
% =====================================================================
\begin{frame}[fragile]{Данные: исполнение в реальном времени}
  \small
  Live-поток от солиста --- последовательность MIDI-событий:
\begin{verbatim}
{"pitch": 64, "timestamp": 1.234, "velocity": 88}
{"pitch": 67, "timestamp": 1.481, "velocity": 95}
...
\end{verbatim}

  \begin{columns}[T,onlytextwidth]
    \column{0.40\textwidth}
    \vspace{0.3em}
    В реальном исполнении \texttt{timestamps} \emph{не совпадают} с
    \texttt{nominal\_onset} из-за rubato, ускорений, замедлений и
    ошибок.

    \vspace{0.5em}
    Это и есть \rlhi{challenge задачи}: трекер должен оставаться
    синхронным, даже когда времени между нотами на лету меняется
    в 2--3 раза.

    \column{0.60\textwidth}
    \centering
    % ГРАФИК 2 — Bar-plot rubato deviations (performance_time − score_onset)
\begin{tikzpicture}
  \begin{axis}[
    width=7.5cm, height=4.5cm,
    ybar,
    bar width=4pt,
    xlabel={\# ноты в партитуре},
    ylabel={$t^{\text{perf}} - t^{\text{nom}}$, c},
    xmin=0, xmax=21,
    ymin=-0.45, ymax=0.45,
    tick style={black},
    axis line style={black},
    grid=major,
    grid style={dotted, gray!50},
    xtick={0,5,10,15,20},
    label style={font=\footnotesize},
    tick label style={font=\scriptsize},
    title style={font=\footnotesize\itshape},
    title={Synthetic rubato: ускорение $\to$ замедление}
  ]
    \addplot[
      fill=black!40,
      draw=black,
    ] coordinates {
      (1, -0.05) (2, -0.12) (3, -0.18) (4, -0.22) (5, -0.28)
      (6, -0.30) (7, -0.28) (8, -0.22) (9, -0.15) (10, -0.08)
      (11, 0.02) (12, 0.10) (13, 0.16) (14, 0.21) (15, 0.27)
      (16, 0.32) (17, 0.36) (18, 0.34) (19, 0.28) (20, 0.20)
    };
    \draw[thick] (axis cs:0,0) -- (axis cs:21,0);
  \end{axis}
\end{tikzpicture}

    \figcaption{Типичный rubato: пианист сначала ускоряется
             (отрицательные дельты), потом замедляется
             (положительные). Это паттерн классического исполнения.}
  \end{columns}
\end{frame}

